% Part: set-theory
% Chapter: story
% Section: grundgesetze

\documentclass[../../../include/open-logic-section]{subfiles}

\begin{document}

\olfileid{sth}{story}{blv}

\olsection{পরিশিষ্ট: ফ্রেগের পঞ্চম মৌলিক বিধি}

\olref[rus]{sec}-এ আমরা ব্যাখ্যা করেছি, রাসেল তাঁর কূটাভাসটিকে
\emph{Grundgesetze}-তে ফ্রেগের বর্ণিত ব্যবস্থার সমস্যা হিসেবে
প্রকাশ করেছিলেন। ফ্রেগের ব্যবস্থায় অনানুষ্ঠানিক ধর্মনির্দেশে
সেট-গঠন সরাসরি বলা ছিল না। তাই এই পরিশিষ্টে খুব সংক্ষেপে
বলব, ফ্রেগের ব্যবস্থায় \emph{কী} ছিল, সেটির সঙ্গে ওই
ধর্মনির্দেশে সেট-গঠনের সম্পর্ক কী এবং রাসেলের কূটাভাসের
সম্পর্কই বা কী।

ফ্রেগের ব্যবস্থা \emph{দ্বিতীয়-ক্রমের}; এটি একটি
\emph{ধারণার বিস্তার} প্রকাশের জন্য তৈরি।\footnote{কঠোরভাবে
বললে, ফ্রেগে আরও সাধারণ একটি ধারণা আনুষ্ঠানিক রূপ দিতে
চেয়েছিলেন: অপেক্ষকের ``মান-বিস্তৃতি''। ধারণার বিস্তার
সেই সাধারণ ধারণার বিশেষ ক্ষেত্র। বিস্তারিত আলোচনার জন্য
\citet[pp.\ 8--9]{Heck2012} দেখুন।} ফ্রেগের অনুপ্রেরণায়
\emph{ধারণা~$F$-এর বিস্তার} বোঝাতে আমরা
$\fregeext{x}{F(x)}$ লিখব। এই রীতি একটি \emph{বিধেয়},
``$F$'', নিয়ে সেটিকে প্রথম-ক্রমের একটি \emph{পদ},
``$\fregeext{x}{F(x)}$'', বানায়। এই রীতির সাহায্যে
ফ্রেগে সদস্যতার নিচের \emph{সংজ্ঞা} দেন:
\[
  a \in b =_\text{df} \exists G(b = \fregeext{x}{G(x)} \land Ga)
\]
মোটামুটি অর্থ: $a \in b$ যদি এবং কেবল যদি $a$ এমন
একটি ধারণার অধীনে পড়ে, যার বিস্তার~$b$। (লক্ষ করুন,
``$\exists G$'' পরিমাণসূচকটি দ্বিতীয়-ক্রমের।) ফ্রেগে
নিচের নীতিটিও গ্রহণ করেন, যার নাম \emph{পঞ্চম মৌলিক বিধি}:
$$\fregeext{x}{F(x)} = \fregeext{x}{G(x)} \liff \forall x (Fx \liff Gx)$$
মোটামুটি অর্থ: দুটি ধারণার বিস্তার একই যদি এবং কেবল যদি
সেগুলি ঠিক একই বস্তুতে প্রযোজ্য হয়। (এখানেও ``$F$'' ও
``$G$''—উভয়ই বিধেয়ের স্থানে আছে।) এখন একটি সরল নীতি
সদস্যতার সঙ্গে ধর্মপরিতৃপ্তির যোগ স্থাপন করে:

\begin{lem}[\emph{Grundgesetze}-তে]\ollabel{lem:Fregeextensions}
$\forall F \forall a(a \in \fregeext{x}{F(x)} \liff Fa)$
\end{lem}

\begin{proof}
$F$ ও~$a$ স্থির করি। এখন $a \in \fregeext{x}{F(x)}$
যদি এবং কেবল যদি $\exists G(\fregeext{x}{F(x)}
= \fregeext{x}{G(x)} \land Ga)$ (সদস্যতার সংজ্ঞা অনুযায়ী);
যদি এবং কেবল যদি $\exists G(\forall x(Fx \liff Gx) \land Ga)$
(পঞ্চম মৌলিক বিধি অনুযায়ী); যদি এবং কেবল যদি $Fa$
(প্রাথমিক দ্বিতীয়-ক্রমের যুক্তিবিদ্যা অনুযায়ী)।
\end{proof}

এ থেকে প্রায় সঙ্গে সঙ্গেই অনানুষ্ঠানিক ধর্মনির্দেশে
সেট-গঠন পাওয়া যায়:

\begin{lem}[\emph{Grundgesetze}-তে]
$\forall F \exists s \forall a (a \in s \liff Fa)$
\end{lem}

\begin{proof}
$F$ স্থির করি। তাহলে \olref{lem:Fregeextensions} থেকে
$\forall a (a \in \fregeext{x}{F(x)} \liff Fa)$ পাওয়া যায়;
তাই অস্তিত্বমূলক সাধারণীকরণে
$\exists s\forall a(a \in s \liff Fa)$। $F$ নির্বিচার
ছিল বলে ফলটি প্রতিষ্ঠিত।
\end{proof}

$F$-কে $\forall x(Fx \liff x \notin x)$ অনুযায়ী নিলেই
রাসেলের কূটাভাস পাওয়া যায়।
% BN-SRC-778 TARGET-CORRECTION-BEGIN
\emph{উৎস-সংশোধন-নোট।} এই শেষ ধাপে নির্দিষ্ট ধর্মটির একটি
বিধেয় আছে বলে ধরতে অনিয়ন্ত্রিত দ্বিতীয়-ক্রমের ধর্মনির্দেশে
গঠন ব্যবহার করা হচ্ছে। পূর্ণ দ্বিতীয়-ক্রমের অর্থতত্ত্বে
বিধেয়-চলের পরিমাণায়নক্ষেত্রে প্রত্যেক বস্তু-উপসেট থাকে; ধর্মটি তাই
পাওয়া যায়। সীমিত বিধেয়ক্ষেত্রের অর্থতত্ত্বে এই গঠন স্বতঃসিদ্ধ
হিসেবে আলাদা করে প্রয়োজন। সেটি বাদ দিয়ে শুধু পঞ্চম মৌলিক
বিধি গ্রহণ করলেই শেষের বিধেয়টির অস্তিত্ব নিশ্চিত হয় না।
% BN-SRC-778 TARGET-CORRECTION-END

\end{document}
